% Part: methods
% Chapter: induction
% Section: inductive-definitions

\documentclass[../../../include/open-logic-section]{subfiles}

\begin{document}

\olfileid{mth}{ind}{idf}

\olsection{ఆగమనాత్మక నిర్వచనాలు}

తర్కశాస్త్రంలో వస్తువుల రకాలను తరచుగా \emph{ఆగమనాత్మకంగా} నిర్వచిస్తాం. అంటే, ఆ రకానికి చెందిన పాత వస్తువుల నుంచి కొత్త వస్తువులను ఎలా పొందాలో తెలిపే నియమాల ద్వారా ఆ రకానికి చెందిన వస్తువులేమిటో నిర్దేశిస్తాం. ఉదాహరణకు, ఒక భాషలోని పదాలు, \tetoken{సూత్రాల}{!!{formula}s} వంటి చిహ్నాల శ్రేణుల ప్రత్యేక రకాలను తరచుగా ఆగమనం ద్వారా నిర్వచిస్తాం. సరళమైన ఉదాహరణగా, $\mathrm{a}$, $\mathrm{b}$, $\mathrm{c}$, $\mathrm{d}$ అక్షరాలు,~$\circ$ చిహ్నం, $[$ మరియు $]$ బ్రాకెట్లతో ఏర్పడిన శ్రేణులను చూడండి. వాటిలో ``$[[\mathrm{c} \circ \mathrm{d}][$'', ``$[\mathrm{a}[]\circ]$'', ``$\mathrm{a}$'', ``$[[\mathrm{a} \circ \mathrm{b}]\circ
\mathrm{d}]$'' ఉన్నాయి. మొదటి రెండు శ్రేణులలో ఏదో ``తప్పు'' ఉందని మీకు అనిపించవచ్చు: మొదటిదానిలో బ్రాకెట్లు అసలు జతపడలేదు; అలాగే ``$\circ$'' అర్థవంతమైన వ్యక్తీకరణలను ``కలపాలి'' అని అనిపించవచ్చు. మూడవ, నాలుగవ శ్రేణులు మెరుగ్గా కనిపిస్తాయి: ప్రతి ``$[$''కు సరిపడే ముగింపు ``$]$'' ఉంది (అటువంటి బ్రాకెట్లు ఉంటే); ప్రతి $\circ$కు రెండు వైపులా బ్రాకెట్ల జంటతో చుట్టబడిన ``చక్కని'' వ్యక్తీకరణలు ఉన్నాయి.

``చక్కని పదం'' అంటే ఏమిటో కచ్చితంగా చెప్పాలనుకుంటున్నాం. మొదట, ప్రతి అక్షరం ఒక్కటే ఉన్నా చక్కని పదమే. కేవలం ఒక్క అక్షరం కానిదేదైనా ``$[t \circ s]$'' రూపంలో ఉండాలి; ఇక్కడ $s$, $t$ స్వయంగా చక్కని పదాలు. తిరిగి, $t$, $s$ చక్కని పదాలైతే వాటి మధ్య $\circ$ ఉంచి, చుట్టూ బ్రాకెట్ల జంటను ఉంచడం ద్వారా కొత్త చక్కని పదాన్ని ఏర్పరచవచ్చు. ఈ క్రియలను ఉపయోగించి చక్కని పదాల సమితిని \emph{నిర్వచించవచ్చు}. ఇదే \emph{ఆగమనాత్మక నిర్వచనం}.

\begin{defn}[చక్కని పదాలు]
  \emph{చక్కని పదాల} సమితిని ఈ విధంగా ఆగమనాత్మకంగా నిర్వచిస్తాం:
  \begin{enumerate}
  \item $\mathrm{a}$, $\mathrm{b}$, $\mathrm{c}$,
    $\mathrm{d}$ అక్షరాలలో ఏదైనా చక్కని పదం.
  \item $s_1$, $s_2$ చక్కని పదాలైతే,
    $[s_1 \circ s_2]$ కూడా చక్కని పదం.
  \item వీటిని మినహాయించి మరేదీ చక్కని పదం కాదు.
  \end{enumerate}
\end{defn}

ఈ నిర్వచనం ప్రకారం, (1),~(2) నిబంధనలతో పరిమిత సంఖ్యలో దశల్లో నిర్మించగలిగితే, అప్పుడూ అప్పుడే ఏదైనా చక్కని పదం అవుతుంది. మొదటి దశలో ఒక్కో అక్షరంతో కూడిన చక్కని పదాలన్నీ నిర్మిస్తాం, అంటే
\[
\mathrm{a}, \mathrm{b}, \mathrm{c}, \mathrm{d}
\]
రెండవ దశలో ఇప్పటికే నిర్మించిన పదాలకు (2)ని వర్తింపజేస్తాం. రెండు అక్షరాల అన్ని జతలకూ
\[
[\mathrm{a} \circ \mathrm{a}], [\mathrm{a} \circ \mathrm{b}],
[\mathrm{b} \circ \mathrm{a}], \dots, [\mathrm{d} \circ \mathrm{d}]
\]
వస్తాయి. మూడవ దశలో, ఇంతవరకు నిర్మించిన ఏ రెండు చక్కని పదాలకైనా (2)ని మళ్లీ వర్తింపజేస్తాం. అప్పుడు $[\mathrm{a} \circ [\mathrm{a} \circ
    \mathrm{a}]]$ వంటి కొత్త చక్కని పదం వస్తుంది---ఇక్కడ $t$ మొదటి దశలోని $\mathrm{a}$, $s$ రెండవ దశలోని $[\mathrm{a} \circ \mathrm{a}]$---అలాగే రెండవ దశలోని $[\mathrm{b} \circ \mathrm{c}]$, $[\mathrm{d}
  \circ \mathrm{b}]$ అనే రెండు పదాలతో $[[\mathrm{b} \circ
    \mathrm{c}] \circ [\mathrm{d} \circ \mathrm{b}]]$ను నిర్మిస్తాం. ఈ విధంగా కొనసాగుతుంది. ఇలా నిర్మించని వస్తువులేవీ చక్కని పదాల సమితిలోకి చేరకుండా (3)వ నిబంధన నిరోధిస్తుంది.

మనకు ``చక్కగా'' అనిపించే ప్రతి చిహ్నశ్రేణీ ఈ నిర్వచనం ప్రకారం చక్కనిదే అని ఇంకా నిరూపించలేదని గమనించండి. అయితే, ఇలా నిర్మించగల ప్రతి శ్రేణీ నిజంగానే ``చక్కగా అనిపిస్తుంది'' అని స్పష్టంగా ఉండాలి: బ్రాకెట్లు జతపడతాయి; $\circ$ స్వయంగా చక్కగా ఉన్న భాగాలను కలుపుతుంది.

ఆగమనాత్మక నిర్వచనాల ముఖ్య లక్షణం ఏమిటంటే, అన్ని చక్కని పదాల గురించి ఏదైనా నిరూపించాలంటే ఏ సందర్భాలను పరిశీలించాలో నిర్వచనమే చెబుతుంది. ఉదాహరణకు, $t$ ఒక చక్కని పదమని తెలిస్తే, $t$ రూపం ఏమై ఉండవచ్చో ఆగమనాత్మక నిర్వచనం చెబుతుంది: $t$ ఒక అక్షరం కావచ్చు; లేదా ఏవో చక్కని పదాలు $s_1$,~$s_2$ కోసం $[s_1 \circ s_2]$ కావచ్చు. (3)వ నిబంధన వల్ల ఇవే సాధ్యాలు.

ఆగమనాత్మకంగా నిర్వచించిన సమితి అంతటినీ గురించిన వాదనలను నిరూపించేటప్పుడు ఆగమనం యొక్క బలమైన రూపం ముఖ్యమవుతుంది. ఉదాహరణకు, పొడవు~$n$ ఉన్న ప్రతి చక్కని పదంలో $[$ల సంఖ్య~$< n/2$ అని నిరూపించాలనుకుందాం. దీన్ని అన్ని~$n$ల గురించిన వాదనగా చూడవచ్చు: ప్రతి $n$కు, పొడవు~$n$ గల ఏ చక్కని పదంలోనైనా $[$ల సంఖ్య $< n/2$.

\begin{prop}
  ఏ $n$కైనా, పొడవు~$n$ ఉన్న చక్కని పదంలో $[$ల సంఖ్య
  $< n/2$.
\end{prop}

\begin{proof}
ఈ ఫలితాన్ని (బలమైన) ఆగమనం ద్వారా నిరూపించడానికి ఈ క్రింది సోపాధిక వాదన సత్యమని చూపాలి:
\begin{quote}
  ప్రతి $l < k$కు, పొడవు~$l$ ఉన్న ఏ చక్కని పదంలోనైనా $[$ల సంఖ్య $< l/2$ అయితే, పొడవు~$k$ ఉన్న ఏ చక్కని పదంలోనైనా $[$ల సంఖ్య $< k/2$.
\end{quote}
దీన్ని చూపడానికి దాని పూర్వపక్షాన్ని తీసుకుందాం: ఏ $l<k$కైనా పొడవు~$l$ గల చక్కని పదాల్లో $[$ల సంఖ్య $< l/2$ అని అనుకుందాం. ఈ పరికల్పనను ఆగమన పరికల్పన అంటాం. పొడవు~$k$ గల చక్కని పదాలకూ ఇదే నిజమని చూపాలి.

అయితే $t$ అనేది పొడవు~$k$ గల చక్కని పదమనుకుందాం. చక్కని పదాలు ఆగమనాత్మకంగా నిర్వచించబడ్డాయి కనుక రెండు సందర్భాలే ఉన్నాయి: (1)~$t$ ఒక్క అక్షరం; లేదా (2)~ఏవో చక్కని పదాలు $s_1$,~$s_2$ కోసం $t$ అనేది $[s_1 \circ s_2]$.
\begin{enumerate}
\item $t$ ఒక అక్షరం. అప్పుడు $k = 1$; $t$లో $[$ల సంఖ్య~$0$. $0 < 1/2$ కనుక వాదన నిలుస్తుంది.
\item ఏవో చక్కని పదాలు $s_1$,~$s_2$ కోసం $t$ అనేది $[s_1 \circ s_2]$. $s_1$ పొడవును $l_1$గా, $s_2$ పొడవును $l_2$గా తీసుకుందాం. అప్పుడు $t$ పొడవు~$k$ అనేది $l_1+l_2+3$ ($s_1$, $s_2$ల పొడవులకు $[$, $\circ$, $]$ అనే మూడు చిహ్నాలు కలిపినది). $l_1+l_2+3$ ఎల్లప్పుడూ $l_1$ కన్నా పెద్దది కనుక $l_1 < k$. అదే విధంగా $l_2
  < k$. అందువల్ల $s_1$,~$s_2$లకు ఆగమన పరికల్పన వర్తిస్తుంది: $s_1$లోని $[$ల సంఖ్య~$m_1$ అనేది $< l_1/2$; $s_2$లోని $[$ల సంఖ్య~$m_2$ అనేది $< l_2/2$.

  $t$లోని $[$ల సంఖ్య, $s_1$లోని $[$ల సంఖ్యకు $s_2$లోని $[$ల సంఖ్య, ఇంకా~$1$ కలిపినది; అంటే $m_1 + m_2 + 1$. $m_1 < l_1/2$, $m_2 < l_2/2$ కనుక:
  \[
  m_1 + m_2 + 1 < \frac{l_1}{2} + \frac{l_2}{2} + 1 = \frac{l_1+l_2+2}{2} < \frac{l_1+l_2+3}{2} = k/2.
  \]
\end{enumerate}
రెండు సందర్భాలలోనూ (ఆగమన పరికల్పన ఆధారంగా) $t$లోని $[$ల సంఖ్య $< k/2$ అని చూపాం. కాబట్టి బలమైన ఆగమనం ద్వారా ప్రతిపాదన వస్తుంది.
\end{proof}

\begin{prob}
  అత్యంత చక్కని పదాల సమితిని ఇలా నిర్వచించండి:
  \begin{enumerate}
  \item $\mathrm{a}$, $\mathrm{b}$, $\mathrm{c}$,
    $\mathrm{d}$ అక్షరాలలో ఏదైనా అత్యంత చక్కని పదం.
  \item $s$ అత్యంత చక్కని పదమైతే $[s]$ కూడా అలాగే ఉంటుంది.
  \item $s_1$, $s_2$ అత్యంత చక్కని పదాలైతే
    $[s_1 \circ s_2]$ కూడా అలాగే ఉంటుంది.
  \item వీటిని మినహాయించి మరేదీ అత్యంత చక్కని పదం కాదు.
  \end{enumerate}
  పొడవు~$n$ గల అత్యంత చక్కని పదం~$t$లో $[$ల సంఖ్య
    $\le n/2 +1$ అని చూపండి.
\end{prob}

\end{document}
